%Root main.tex
%---------------------------------------------------------------------------
%付録A　付録
%---------------------------------------------------------------------------

%\lstnewenvironment{mylisting}[1][]
%    {\lstset{
%        frame=single,
%        numbersep=10pt,
%        tabsize=2,
%        #1
%    }
%}{}
\begin{mylisting}[language=c,caption=SVPWMモジュール,label=hurokuB_SVPWM]

#include <Stdlib.h>
#include <String.h>
#include <math.h>
#include <Psim.h>

// PLACE GLOBAL VARIABLES OR USER FUNCTIONS HERE...
int n = 5;
//double E = 116;//1.18/100;;//116;
//OPEN or PI
//double E = 200/(sqrt(3));//1.18/100;;//116;
//MSSRDBDC
double E = 2.0e-4;//1.1547e-4;//0.10235;//1.1547e-4;//0.0000036084;//
int cnt=0;
int cnt_max = 200;

double Ts,Tc;
double fs,fc;

double ua,ub;
double sector,type;
double theta_deg,theta,theta_PI;
double ma,ma_l,ma_buf;
double ua_1,ub_1,theta_1;
int H,L,R;
double t1,t2,t0;
int UP,UN,VP,VN,WP,WN;
double T0,T1,T2,T3,T4,T5,T6,Tlevel;
double T1_buf,T2_buf,T3_buf,T4_buf,T5_buf,T6_buf;
double svm_t,Tcnt;
double SU,SV,SW;


/////////////////////////////////////////////////////////////////////
// FUNCTION: SimulationStep
//   This function runs at every time step.
//double t: (read only) time
//double delt: (read only) time step as in Simulation control
//double *in: (read only) zero based array of input values. in[0] is the first node, in[1] second input...
//double *out: (write only) zero based array of output values. out[0] is the first node, out[1] second output...
//int *pnError: (write only)  assign  *pnError = 1;  if there is an error and set the error message in szErrorMsg
//    strcpy(szErrorMsg, "Error message here..."); 
// DO NOT CHANGE THE NAME OR PARAMETERS OF THIS FUNCTION
void SimulationStep(
		double t, double delt, double *in, double *out,
		 int *pnError, char * szErrorMsg,
		 void ** reserved_UserData, int reserved_ThreadIndex, void * reserved_AppPtr)
{
// ENTER YOUR CODE HERE...
//入力
ua = in[0];
ub = in[1];
Ts = in[2];
Tc = in[3];

Tcnt = Tc/cnt_max;

//非MSの場合はカウンタ = 0の時のみ演算する
if(cnt == 0.0){

/***操作量の比較によりセクタ導出****/
if(ua>=0&&ub>=0)
{if(sqrt(3)*ua>=ub){sector = 1;	}else{sector = 2;}
}else if(ua<=0&&ub>=0){if(-sqrt(3)*ua>ub){sector = 3;}else{sector = 2;}
}else if(ua<=0&&ub<=0){if(-sqrt(3)*ua>=-ub){sector = 4;}else{sector = 5;}
}else if(ua>=0&&ub<=0){if(sqrt(3)*ua<=-ub){sector = 5;}else{sector = 6;}
}else{sector = 0;}

}//cnt_if

/***操作量から角度導出(MS)****/
	theta = atan2(ub,ua);
	if(theta >= 0){
	theta_PI = theta;
	theta_deg = (180.0 / M_PI) * theta_PI;
}else{
		theta_PI = 2*M_PI + theta;
		theta_deg = (180.0 / M_PI) * theta_PI;
	}


/***ma導出(MS)****/
ma_l =sqrt(ua*ua+ub*ub);
ma_buf =ma_l/E;
if(ma_buf >=1.0){ma = 1.0;}
else{ma = ma_buf;}

/***sector1に遷移****/
if(sector == 1){
	ua_1 = ua;
	ub_1 = ub;
}else if(sector == 2){
	ua_1 = ma_l*cos(theta - M_PI/3);
	ub_1 = ma_l*sin(theta - M_PI/3);
}else if(sector == 3){
	ua_1 = ma_l*cos(theta - 2*M_PI/3);
	ub_1 = ma_l*sin(theta - 2*M_PI/3);
}else if(sector == 4){
	ua_1 = ma_l*cos(theta - M_PI);
	ub_1 = ma_l*sin(theta - M_PI);
}else if(sector == 5){
	ua_1 = ma_l*cos(theta - 4*M_PI/3);
	ub_1 = ma_l*sin(theta - 4*M_PI/3);
}else if(sector == 6){
	ua_1 = ma_l*cos(theta - 5*M_PI/3);
	ub_1 = ma_l*sin(theta - 5*M_PI/3);
}
theta_1 = atan2(ub_1,ua_1);

if(cnt == 0.0){//anti ms
/***H,L,R導出****/
H = abs((n-1)*ma*cos(M_PI/6 - theta_1))+1;
L = abs((n-1)*ma*sin(theta_1))+1;
R = abs((n-1)*ma*sin((M_PI/3)-theta_1))+1;

/***secor type導出****/
if(H == 1){
	if(sector == 1|| sector == 3 || sector == 5){type = 1;
	}else if (sector == 2|| sector == 4 || sector == 6){type = 2;
	}else{type = 0;}
}else if(H == 2){
	if(L == 1){if(R == 2){type = 1;}else{type = 2;}
	}else if(L == 2){type = 1;}
}else if(H == 3){
	if(L == 1){if(R == 3){type = 1;}else{type = 2;}
	}else if(L == 2){if(R == 2){type = 1;}else{type = 2;}
	}else if(L == 3){type = 1;}}
else if(H == 4){
	if(L == 1){if(R == 4){type = 1;}else{type = 2;}
	}else if(L == 2){if(R == 3){type = 1;}else{type = 2;}
	}else if(L == 3){if(R == 2){type = 1;}else{type = 2;}
	}else if(L == 4){type = 1;}
}

}//cnt_if

/***t1,t2導出****/
if(type == 1){
t1 = Tc*((n-1) * ma * sin((M_PI/3)-theta_1) - (H-L));
t2 = Tc*((n-1)*ma*sin(theta_1)-(L-1));
}else if(type == 2){
t1 = Tc*((H-L)-(n-1)*ma*sin((M_PI/3)-theta_1));
t2 = Tc*((n-1)*ma*sin((M_PI/3)+theta_1)-(H-1));
}else{t1 = 0;t2 = 0;}
t0 = Tc - t2 -t1;

/***スイッチングパルス導出****/
//if(cnt == 0.0){
		T1 = t0/4;
		T2 = t0/4 + t1/2;
		T3 = t0/4 + t1/2 + t2/2;
		T4 = t0/4 + t1/2 + t2/2 + t0/2;
		T5 = t0/4 + t1/2 + t2 + t0/2;
		T6 = t0/4 + t1 + t2 + t0/2;
//	}//if(cnt == 0.0){はここで終了
	if(svm_t < T1){UP = 0,VP = 0,WP = 0,UN = 1,VN = 1,WN = 1,Tlevel = 0;}
	else if((T1 < svm_t) && (svm_t <= T2)){UP = 1,VP = 0,WP = 0,UN = 0,VN = 1,WN = 1,Tlevel = 1;}
	else if((T2 < svm_t) && (svm_t <= T3)){UP = 1,VP = 1,WP = 0,UN = 0,VN = 0,WN = 1,Tlevel = 2;}
	else if((T3 < svm_t) && (svm_t <= T4)){UP = 1,VP = 1,WP = 1,UN = 0,VN = 0,WN = 0,Tlevel = 3;}
	else if((T4 < svm_t) && (svm_t <= T5)){UP = 1,VP = 1,WP = 0,UN = 0,VN = 0,WN = 1,Tlevel = 4;}
	else if((T5 < svm_t) && (svm_t <= T6)){UP = 1,VP = 0,WP = 0,UN = 0,VN = 0,WN = 1,Tlevel = 5;}
	else if((T6 < svm_t) && (svm_t <= Tc)){UP = 0,VP = 0,WP = 0,UN = 1,VN = 1,WN = 1,Tlevel = 6;}

/***スイッチングパルス導出(W相ON時間0)****/
/*if(cnt == 0.0){
		T1 = t0/2;
		T2 = t0/2 + t1/2;
		T3 = t0/2 + t1/2 + t2/2;
		T4 = t0/2 + t1/2 + t2;
		T5 = t0/2 + t1 + t2;
	}
	if(svm_t < T1){UP = 0,VP = 0,WP = 0,UN = 1,VN = 1,WN = 1,Tlevel = 0;}
	else if((T1 < svm_t) && (svm_t <= T2)){UP = 1,VP = 0,WP = 0,UN = 0,VN = 1,WN = 1,Tlevel = 1;}
	else if((T2 < svm_t) && (svm_t <= T3)){UP = 1,VP = 1,WP = 0,UN = 0,VN = 0,WN = 1,Tlevel = 2;}
	else if((T3 < svm_t) && (svm_t <= T4)){UP = 1,VP = 1,WP = 0,UN = 0,VN = 0,WN = 1,Tlevel = 4;}
	else if((T4 < svm_t) && (svm_t <= T5)){UP = 1,VP = 0,WP = 0,UN = 0,VN = 0,WN = 1,Tlevel = 5;}
	else if((T5 < svm_t) && (svm_t <= Tc)){UP = 0,VP = 0,WP = 0,UN = 1,VN = 1,WN = 1,Tlevel = 6;}*/
else{UP = 0,VP = 0,WP = 0,UN = 0,VN = 0,WN = 0,Tlevel = -1;}

/***SU,SV,SW導出（スイッチングレベル付与）****/
if(sector == 1){
	if(type == 1){
		SU = UP + (H-1);
		SV = VP + (L-1);
		SW = WP;
		}else if(type == 2){
		SU = VP +(H-1);
		SV = UP +(L-1);
		SW = WP;
		}
}else if(sector == 2){
	if(type == 1){
		SV = UP + (H-1);
		SU = VP + (R-1);
		SW = WP;
		}else if(type == 2){
		SV = VP +(H-1);
		SU = UP +(R-1);
		SW = WP;
		}
}else if(sector == 3){
	if(type == 1){
		SV = UP + (H-1);
		SW = VP + (L-1);
		SU = WP;
		}else if(type == 2){
		SV = VP +(H-1);
		SW = UP +(L-1);
		SU = WP;
		}
}else if(sector == 4){
	if(type == 1){
		SW = UP + (H-1);
		SV = VP + (R-1);
		SU = WP;
		}else if(type == 2){
		SW = VP +(H-1);
		SV = UP +(R-1);
		SU = WP;
		}
}else if(sector == 5){
	if(type == 1){
		SW = UP + (H-1);
		SU = VP + (L-1);
		SV = WP;
		}else if(type == 2){
		SW = VP +(H-1);
		SU = UP +(L-1);
		SV = WP;
		}
}else if(sector == 6){
	if(type ==1){
		SU = UP + (H-1);
		SW = VP + (R-1);
		SV = WP;
		}else if(type == 2){
		SU = VP +(H-1);
		SW = UP +(R-1);
		SV = WP;
		}
}

//out
out[0] = sector;
out[1] = theta_PI;
out[2] = theta_deg;
out[3] = ma_l;
out[4] = ma;
out[5] = ua_1;
out[6] = ub_1;
out[7] = theta_1;
out[8] = H;
out[9] = L;
out[10] = R;
out[11] = type;
out[12] = t1;                  
out[13] = t2;
out[14] = t0;
out[15] = UP;
out[16] = UN;
out[17] = VP;
out[18] = VN;
out[19] = WP;
out[20] = WN;
out[21] = Tlevel;
out[22] = cnt;
out[23] = svm_t;
out[24] = SU;
out[25] = SV;
out[26] = SW;


if(cnt == cnt_max){
	cnt = 0;
	svm_t = 0.0;
}else{
	cnt = cnt + 1;
	svm_t += Tcnt; 
}

}


/////////////////////////////////////////////////////////////////////
// FUNCTION: SimulationBegin
//   Initialization function. This function runs once at the beginning of simulation
//   For parameter sweep or AC sweep simulation, this function runs at the beginning of each simulation cycle.
//   Use this function to initialize static or global variables.
//const char *szId: (read only) Name of the C-block 
//int nInputCount: (read only) Number of input nodes
//int nOutputCount: (read only) Number of output nodes
//int nParameterCount: (read only) Number of parameters is always zero for C-Blocks.  Ignore nParameterCount and pszParameters
//int *pnError: (write only)  assign  *pnError = 1;  if there is an error and set the error message in szErrorMsg
//    strcpy(szErrorMsg, "Error message here..."); 
// DO NOT CHANGE THE NAME OR PARAMETERS OF THIS FUNCTION
void SimulationBegin(
		const char *szId, int nInputCount, int nOutputCount,
		 int nParameterCount, const char ** pszParameters,
		 int *pnError, char * szErrorMsg,
		 void ** reserved_UserData, int reserved_ThreadIndex, void * reserved_AppPtr)
{
// ENTER INITIALIZATION CODE HERE...




}


/////////////////////////////////////////////////////////////////////
// FUNCTION: SimulationEnd
//   Termination function. This function runs once at the end of simulation
//   For parameter sweep or AC sweep simulation, this function runs at the end of each simulation cycle.
//   Use this function to de-allocate any allocated memory or to save the result of simulation in an alternate file.
// Ignore all parameters for C-block 
// DO NOT CHANGE THE NAME OR PARAMETERS OF THIS FUNCTION
void SimulationEnd(const char *szId, void ** reserved_UserData, int reserved_ThreadIndex, void * reserved_AppPtr)
{




}





\end{mylisting}

